%\chapter{Тест}
% TODO: Написати Розділ 1
\section{Аналіз предметної області}

\subsection{Загальні положення}

\subsection{Огляд існуючих рішень}

\subsubsection{gRPC}

%Сценарії використання протоколу JSTP його користувачем (розробником прикладних
%додатків з використанням протоколу) наведено у
%таблицях~\ref{tbl:usecase1}--\ref{tbl:usecase18}.

%\begin{table}[htbp]
%\label{tbl:usecase1}
%  \captionbox{юк1}{%
%    \begin{tabu} to \textwidth { | X[1,l] | X[4,l] | }
%      \hline
%      Назва & Створення віддаленого застосунку \\
%      \hline
%      Опис & Користувач має можливість створити віддалений застосунок, вказавши
%      його ім'я, список інтерфейсів і методів в них, та список реакцій на
%      події, що надсилає клієнт \\
%      \hline
%      Учасники & Користувач \\
%      \hline
%      Передумови & -- \\
%      \hline
%      Постумови & Користувач створив застосунок \\
%      \hline
%    \end{tabu}
%  }
%\end{table}
%
%\begin{table}[htbp]
%\label{tbl:usecase2}
%  \captionbox{юк2}{%
%    \begin{tabu} to \textwidth { | X[1,l] | X[4,l] | }
%      \hline
%      Назва & Створення серверу \\
%      \hline
%      Опис & Користувач має можливість створити сервер, що буде обробляти запити до віддаленого застосунку, для цього необхідно вказати мережевий транспорт, адресу серверу та список віддалених застосунків, що буде ним обслуговуватись \\
%      \hline
%      Учасники & Користувач \\
%      \hline
%      Передумови & Користувач створив застосунок \\
%      \hline
%      Постумови & Користувач створив сервер \\
%      \hline
%    \end{tabu}
%  }
%\end{table}

%\begin{table}[htbp]
%  \captionbox{\label{tbl:usecase3}}{%
%    \begin{tabu} to \textwidth { | X[1,l] | X[4,l] | }
%      \hline
%      Назва & Створення серверу з різними версіями віддаленого застосунку \\
%      \hline
%      Опис & Користувач має можливість створити декілька версій віддаленого застосунку, вказавши його ім’я, версії, список інтерфейсів і методів в них, та список реакцій на події, що надсилає клієнт \\
%      \hline
%      Учасники & Користувач \\
%      \hline
%      Передумови & --- \\
%      \hline
%      Постумови & Користувач сервер з різними версіями застосунку \\
%      \hline
%    \end{tabu}
%  }
%\end{table}

%\begin{table}[htbp]
%  \captionbox{\label{tbl:usecase4}}{%
%    \begin{tabu} to \textwidth { | X[1,l] | X[4,l] | }
%      \hline
%      Назва & Створення клієнту \\
%      \hline
%      Опис & Користувач має можливість створити клієнт, що буде обробляти запити до віддаленого застосунку, для цього необхідно вказати застосунок, що буде доступний через мережеве підключення \\
%      \hline
%      Учасники & Користувач \\
%      \hline
%      Передумови & Користувач створив застосунок \\
%      \hline
%      Постумови & Користувач створив клієнт \\
%      \hline
%    \end{tabu}
%  }
%\end{table}
%
%\begin{table}[htbp]
%  \captionbox{\label{tbl:usecase5}}{%
%    \begin{tabu} to \textwidth { | X[1,l] | X[4,l] | }
%      \hline
%      Назва & Підключення до серверу \\
%      \hline
%      Опис & Користувач має можливість створити підключення між клієнтом та сервером, для цього необхідно вказати транспорт й адресу серверу, назву та версію додатку, до якого буде здійснено підключення \\
%      \hline
%      Учасники & Користувач \\
%      \hline
%      Передумови & Користувач створив сервер \\
%      \hline
%      Постумови & Користувач встановив підключення між клієнтом та сервером \\
%      \hline
%    \end{tabu}
%  }
%\end{table}
%
%\begin{table}[htbp]
%  \captionbox{\label{tbl:usecase6}}{%
%    \begin{tabu} to \textwidth { | X[1,l] | X[4,l] | }
%      \hline
%      Назва & Підключення до вказаної версії віддаленого застосунку \\
%      \hline
%      Опис & Користувач має можливість вказати точну версію або проміжок, якому повинна відповідати версія віддаленого застосунку, клієнта буде підключено до найновішої підходящої версії \\
%      \hline
%      Учасники & Користувач \\
%      \hline
%      Передумови & Користувач створив сервер \\
%      \hline
%      Постумови & Клієнт підключено до найновішої підходящої версії віддаленого застосунку \\
%      \hline
%    \end{tabu}
%  }
%\end{table}
%
%\begin{table}[htbp]
%  \captionbox{\label{tbl:usecase7}}{%
%    \begin{tabu} to \textwidth { | X[1,l] | X[4,l] | }
%      \hline
%      Назва & Автентифікація \\
%      \hline
%      Опис & Користувач має можливість автентифікувати клієнт під час підключення до серверу, для цього необхідно створити політику автентифікації, при створенні серверу та передати правильні дані для автентифікації під час підключення \\
%      \hline
%      Учасники & Користувач \\
%      \hline
%      Передумови & Користувач створив сервер з політикою автентифікації \\
%      \hline
%      Постумови & Клієнт було автентифіковано \\
%      \hline
%    \end{tabu}
%  }
%\end{table}
%
%\begin{table}[htbp]
%  \captionbox{\label{tbl:usecase8}}{%
%    \begin{tabu} to \textwidth { | X[1,l] | X[4,l] | }
%      \hline
%      Назва & Відключення \\
%      \hline
%      Опис & Користувач має можливість відключитись від сервера, пославши відповідне повідомлення \\
%      \hline
%      Учасники & Користувач \\
%      \hline
%      Передумови & Клієнт було підключено до сервера \\
%      \hline
%      Постумови & Клієнт було відключено від сервера \\
%      \hline
%    \end{tabu}
%  }
%\end{table}
%
%\begin{table}[htbp]
%  \captionbox{\label{tbl:usecase9}}{%
%    \begin{tabu} to \textwidth { | X[1,l] | X[4,l] | }
%      \hline
%      Назва & Відновлення підключення \\
%      \hline
%      Опис & Користувач має можливість повторно підключитись до існуючої сесії \\
%      \hline
%      Учасники & Користувач \\
%      \hline
%      Передумови & На сервері існує сесія \\
%      \hline
%      Постумови & Клієнт було перепідключено до існуючої сесії \\
%      \hline
%    \end{tabu}
%  }
%\end{table}
%
%\begin{table}[htbp]
%  \captionbox{\label{tbl:usecase10}}{%
%    \begin{tabu} to \textwidth { | X[1,l] | X[4,l] | }
%      \hline
%      Назва & Відновлення підключення при розриві мережевого підключення \\
%      \hline
%      Опис & Клієнт автоматично відновлює сесію, якщо підключення було перервано в зв’язку з мережевою помилкою \\
%      \hline
%      Учасники & Користувач \\
%      \hline
%      Передумови & На сервері існує сесія, розрив мережевого підключення \\
%      \hline
%      Постумови & Клієнт було автоматично перепідключено до існуючої сесії \\
%      \hline
%    \end{tabu}
%  }
%\end{table}
%
%\begin{table}[htbp]
%  \captionbox{\label{tbl:usecase11}}{%
%    \begin{tabu} to \textwidth { | X[1,l] | X[4,l] | }
%      \hline
%      Назва & Виконати ідемпотентний виклик віддаленої процедури \\
%      \hline
%      Опис & Користувач має можливість здійснити ідемпотентний виклик віддаленої процедури, для цього одна із сторін посилає відповідний запит, інша обробляє його та повертає результат, якщо відбулась мережева помилка, віддалену функцію буде викликано повторно \\
%      \hline
%      Учасники & Користувач \\
%      \hline
%      Передумови & Клієнт підключено до серверу \\
%      \hline
%      Постумови & Результат виконання віддаленої процедури було надіслано стороні-ініціатору \\
%      \hline
%    \end{tabu}
%  }
%\end{table}
%
%\begin{table}[htbp]
%  \captionbox{\label{tbl:usecase12}}{%
%    \begin{tabu} to \textwidth { | X[1,l] | X[4,l] | }
%      \hline
%      Назва & Виконати не ідемпотентний виклик віддаленої процедури \\
%      \hline
%      Опис & Користувач має можливість здійснити не ідемпотентний виклик віддаленої процедури, для цього одна із сторін посилає відповідний запит, інша обробляє його та повертає результат, якщо відбулась мережева помилка, віддалену функцію не буде викликано повторно \\
%      \hline
%      Учасники & Користувач \\
%      \hline
%      Передумови & Клієнт підключено до серверу \\
%      \hline
%      Постумови & Результат виконання віддаленої процедури було надіслано стороні-ініціатору \\
%      \hline
%    \end{tabu}
%  }
%\end{table}
%
%\begin{table}[htbp]
%  \captionbox{\label{tbl:usecase13}}{%
%    \begin{tabu} to \textwidth { | X[1,l] | X[4,l] | }
%      \hline
%      Назва & Генерація події \\
%      \hline
%      Опис & Користувач має можливість згенерувати подію на одній із сторін, що буде передана іншій \\
%      \hline
%      Учасники & Користувач \\
%      \hline
%      Передумови & Клієнт підключено до серверу \\
%      \hline
%      Постумови & Повідомлення про подію було отримано іншою стороною \\
%      \hline
%    \end{tabu}
%  }
%\end{table}
%
%\begin{table}[htbp]
%  \captionbox{\label{tbl:usecase14}}{%
%    \begin{tabu} to \textwidth { | X[1,l] | X[4,l] | }
%      \hline
%      Назва & Підписки на подію \\
%      \hline
%      Опис & Користувач має можливість підписати функцію на певну подію, ця функція буде викликана у разі генерації події іншою стороною \\
%      \hline
%      Учасники & Користувач \\
%      \hline
%      Передумови & Клієнт підключено до серверу \\
%      \hline
%      Постумови & Одна зі сторін була підписана на подію \\
%      \hline
%    \end{tabu}
%  }
%\end{table}
%
%\begin{table}[htbp]
%  \captionbox{\label{tbl:usecase15}}{%
%    \begin{tabu} to \textwidth { | X[1,l] | X[4,l] | }
%      \hline
%      Назва & Обробка подій \\
%      \hline
%      Опис & Користувач має можливість реагувати на подію, згенеровані іншою стороною підключення \\
%      \hline
%      Учасники & Користувач \\
%      \hline
%      Передумови & Одна зі сторін підключення підписується на подію, інша генерує її \\
%      \hline
%      Постумови & Подія була оброблена \\
%      \hline
%    \end{tabu}
%  }
%\end{table}
%
%\begin{table}[htbp]
%  \captionbox{\label{tbl:usecase16}}{%
%    \begin{tabu} to \textwidth { | X[1,l] | X[4,l] | }
%      \hline
%      Назва & Відміна підписки на подію \\
%      \hline
%      Опис & Користувач має можливість відмінити підписку на подію \\
%      \hline
%      Учасники & Користувач \\
%      \hline
%      Передумови & Одна зі сторін підключення підписана на подію \\
%      \hline
%      Постумови & Підписку на подію було відмінено \\
%      \hline
%    \end{tabu}
%  }
%\end{table}
%
%\begin{table}[htbp]
%  \captionbox{\label{tbl:usecase17}}{%
%    \begin{tabu} to \textwidth { | X[1,l] | X[4,l] | }
%      \hline
%      Назва & Старт heartbeat \\
%      \hline
%      Опис & Користувач має можливість ініціювати періодичну відправку ping повідомлень \\
%      \hline
%      Учасники & Користувач \\
%      \hline
%      Передумови & Клієнт підключений до сервера \\
%      \hline
%      Постумови & Періодична відправка ping повідомлень була успішно ініційована \\
%      \hline
%    \end{tabu}
%  }
%\end{table}
%
%\begin{table}[htbp]
%  \captionbox{\label{tbl:usecase18}}{%
%    \begin{tabu} to \textwidth { | X[1,l] | X[4,l] | }
%      \hline
%      Назва & Стоп heartbeat \\
%      \hline
%      Опис & Користувач має можливість зупинити періодичну відправку ping повідомлень \\
%      \hline
%      Учасники & Користувач \\
%      \hline
%      Передумови & Періодична відправка ping повідомлень активована \\
%      \hline
%      Постумови & Періодична відправка ping повідомлень була успішно зупинена \\
%      \hline
%    \end{tabu}
%  }
%\end{table}

\begin{unnumbered}

\section{Висновки до розділу 1}

\end{unnumbered}
